\documentclass[twocolumn,pra,amsmath,amssymb,superscriptaddress]{revtex4-2}

\usepackage{graphicx}
\usepackage{hyperref}
\usepackage{booktabs}
\usepackage{xcolor}
\usepackage{amsthm}

\newtheorem{theorem}{Theorem}
\newtheorem{proposition}[theorem]{Proposition}
\newtheorem{corollary}[theorem]{Corollary}
\newtheorem{lemma}[theorem]{Lemma}
\newtheorem{remark}[theorem]{Remark}

\begin{document}

\title{Two-body interactions from the tensor-product spectral triple:\\
{\large Selection rules from the algebra, coupling strengths from the metric}}

\author{J.~Loutey}
\affiliation{Independent Researcher, Kent, Washington}

\date{June 1, 2026}

\begin{abstract}
We construct the tensor-product spectral triple
$\mathcal{T}_{S^3} \otimes \mathcal{T}_{S^3}$ for two particles on the
Fock-projected three-sphere at finite cutoff $n_{\max}$, and investigate
whether the gauged spectral action generates the Coulomb inter-particle
interaction from the framework's own algebraic structure. The free
tensor-product Dirac $D_{\mathrm{total}} = D_1 \otimes I + \gamma_1
\otimes D_2$ satisfies $\{D, \gamma\} = 0$ exactly, so
$D_{\mathrm{total}}^2$ factorizes and the free spectral action contains
no interaction (Theorem~\ref{thm:free_factorization}). Gauging via the
inner fluctuation $A = \sum_i a_i^\dagger [D, a_i]$ from the truncated
operator system $\mathcal{O}_{n_{\max}}$ produces a genuine two-body
interaction: the connected (non-factorizable) fraction of $\{D, A\}$ is
$77\%$ at $n_{\max} = 2$ and $32\%$ at $n_{\max} = 3$
(Proposition~\ref{prop:connected}). The angular decomposition of this
interaction is $100\%$ $m$-conserving, $100\%$ Gaunt-compatible, and
pure $k = 0$ monopole at both cutoffs
(Theorem~\ref{thm:angular_structure}). However, the radial matrix
element ratios do not match the Coulomb $1/r_{12}$ kernel: the
coefficient of variation exceeds $1600\%$ and the Pearson correlation
is $0.41$--$0.58$ (Proposition~\ref{prop:radial_negative}). These
results precisely delineate Paper~31's universal/Coulomb partition at
the two-body level: the \emph{algebra} determines which channels
couple (angular selection rules); the \emph{metric} determines how
strongly (radial coupling strengths). The multi-focal composition wall
identified across six independent observables in the GeoVac programme
is structurally located at this algebra/metric boundary.
\end{abstract}

\maketitle

%% =====================================================================
\section{Introduction}
\label{sec:intro}
%% =====================================================================

The GeoVac framework constructs a discrete spectral triple
$(\mathcal{A}, \mathcal{H}, D)$ on the Fock-projected
three-sphere~\cite{paper7,paper32}: the algebra $\mathcal{A}$ of
truncated multipliers on $S^3$, the Hilbert space $\mathcal{H}$ of
Fock-projected states, and the Camporesi--Higuchi Dirac operator $D$.
This single-particle spectral triple has been shown to converge to the
continuum $S^3$ in the van~Suijlekom state-space
Gromov--Hausdorff distance~\cite{paper38}, and the spectral action
$\mathrm{Tr}\,f(D^2/\Lambda^2)$ reproduces the correct gravity sector
at one loop~\cite{paper51}.

A central open question is whether the framework can derive
\emph{inter-particle interactions} from its own structure, rather than
importing them as external input. The electron--electron Coulomb
repulsion $V_{ee} = 1/r_{12}$, the nucleon--nucleon force, and the
core--valence interaction in molecular systems are all two-body
couplings that the framework currently imports from standard quantum
mechanics. The ``multi-focal composition wall''~\cite{paper34} ---
the observation that the framework couples discrete quantum numbers
cleanly but has no native composition theorem for spatial focal
lengths --- has been documented across six independent
observables~\cite{CLAUDE_wall}.

In this paper we attack the problem directly. We construct the
tensor-product spectral triple $\mathcal{T}_1 \otimes \mathcal{T}_2$
for two particles, each on a copy of the truncated $S^3$, and compute
the gauged spectral action to determine what interaction, if any,
emerges from the algebra's inner fluctuations.

The results are sharp: the algebra produces the correct angular
selection rules (Gaunt compatibility, $m$-conservation, monopole
dominance), but the radial coupling strengths are not proportional to
the Coulomb kernel. This places the multi-focal wall precisely at the
boundary between the algebra $\mathcal{A}$ (which determines angular
structure) and the Dirac operator $D$ (which determines radial
couplings) --- the universal/Coulomb partition of
Paper~31~\cite{paper31}, now verified at the two-body level.

%% =====================================================================
\section{Setup}
\label{sec:setup}
%% =====================================================================

\subsection{Single-particle spectral triple}

We work with the scalar sector of the Fock-projected $S^3$ at cutoff
$n_{\max}$. The Hilbert space is
\begin{equation}
  \mathcal{H}_{n_{\max}} =
  \mathrm{span}\{|n, l, m\rangle : 1 \le n \le n_{\max}\},
  \label{eq:H_scalar}
\end{equation}
with $\dim \mathcal{H}_{n_{\max}} = N = \sum_{n=1}^{n_{\max}} n^2$.
For the Dirac structure, we chirality-double to $\hat{\mathcal{H}} =
\mathcal{H}^+ \oplus \mathcal{H}^-$ with $\dim \hat{\mathcal{H}} =
2N$.

The Dirac operator is the Camporesi--Higuchi form~\cite{camporesi_higuchi1996}:
\begin{equation}
  D |n, l, m, \chi\rangle = \chi \left(n + \tfrac{1}{2}\right)
  |n, l, m, \chi\rangle,
  \label{eq:dirac}
\end{equation}
where $\chi = \pm 1$ is the chirality. The chirality operator
$\gamma$ swaps the two sectors: $\gamma |n, l, m, +\rangle = |n, l,
m, -\rangle$.

The algebra $\mathcal{A} = \mathcal{O}_{n_{\max}}$ is the Connes--van
Suijlekom~\cite{connes_vs2021} truncated operator system, generated
by the $3$-$Y$ multiplier matrices $M_{NLM}$ with entries
\begin{equation}
  (M_{NLM})_{(nlm),(n'l'm')} =
  \int_{S^3} \overline{Y^{(3)}_{nlm}} \cdot Y^{(3)}_{NLM}
  \cdot Y^{(3)}_{n'l'm'} \, d\Omega.
  \label{eq:multiplier}
\end{equation}
These are embedded into $\hat{\mathcal{H}}$ as $\hat{M} =
\mathrm{diag}(M, M)$, acting identically on both chirality sectors.

\subsection{Tensor-product Dirac}

For two particles, the total Hilbert space is $\hat{\mathcal{H}}_1
\otimes \hat{\mathcal{H}}_2$ with dimension $(2N)^2$. The
tensor-product Dirac operator is the standard
product~\cite{connes_book1994}:
\begin{equation}
  D_{\mathrm{total}} = D_1 \otimes I_2 + \gamma_1 \otimes D_2.
  \label{eq:D_total}
\end{equation}
The chirality twist $\gamma_1$ ensures $D_{\mathrm{total}}$ is a
first-order operator on the graded tensor product.

%% =====================================================================
\section{Free factorization}
\label{sec:free}
%% =====================================================================

\begin{theorem}[Free factorization]
\label{thm:free_factorization}
For the Camporesi--Higuchi Dirac~\eqref{eq:dirac} on the
chirality-doubled scalar basis, $\{D, \gamma\} = 0$ exactly. As a
consequence,
\begin{equation}
  D_{\mathrm{total}}^2 = D_1^2 \otimes I_2 + I_1 \otimes D_2^2,
  \label{eq:D2_factorizes}
\end{equation}
and the spectral action $\mathrm{Tr}\,f(D_{\mathrm{total}}^2/\Lambda^2)
= \mathrm{Tr}\,f(D_1^2/\Lambda^2) \cdot \mathrm{Tr}\,
f(D_2^2/\Lambda^2)$ factorizes exactly at every cutoff $\Lambda$.
\end{theorem}

\begin{proof}
The Dirac eigenvalue on $|n, l, m, \chi\rangle$ is $\chi(n +
\frac{1}{2})$. The chirality operator sends $\chi \to -\chi$. For any
basis state,
\[
(D\gamma + \gamma D)|n,l,m,\chi\rangle =
\chi(n+\tfrac{1}{2})(-\chi)(n+\tfrac{1}{2})
+ (-\chi)(n+\tfrac{1}{2})\chi(n+\tfrac{1}{2}) = 0,
\]
using $D|n,l,m,-\chi\rangle = -\chi(n+\frac{1}{2})|n,l,m,-\chi\rangle$.
Hence $\{D,\gamma\} = 0$. Expanding $D_{\mathrm{total}}^2$:
\[
D_{\mathrm{total}}^2 = D_1^2 \otimes I + (D_1 \gamma_1 + \gamma_1
D_1) \otimes D_2 + I \otimes D_2^2.
\]
The cross term vanishes by $\{D_1, \gamma_1\} = 0$.
\end{proof}

\begin{remark}
This is the correct physical statement: free particles do not
interact. The spectral action of the free Dirac gives independent
single-particle kinetic energies. Any interaction must enter through
\emph{gauging}.
\end{remark}

\emph{Numerical verification.} At $n_{\max} \in \{2, 3\}$, the heat
trace ratio $\mathrm{Tr}(e^{-t D_{\mathrm{total}}^2}) /
[\mathrm{Tr}(e^{-t D_1^2}) \cdot \mathrm{Tr}(e^{-t D_2^2})] =
1.000000$ to machine precision at all tested $t \in \{0.1, 0.5, 1.0,
2.0\}$.

%% =====================================================================
\section{Gauged interaction}
\label{sec:gauged}
%% =====================================================================

The inner fluctuation of $D_{\mathrm{total}}$ by the tensor-product
algebra $\mathcal{A}_1 \otimes \mathcal{A}_2$ is
\begin{equation}
  [D_{\mathrm{total}}, \hat{M}_i \otimes \hat{M}_j] =
  [D_1, \hat{M}_i] \otimes \hat{M}_j
  + \gamma_1 \hat{M}_i \otimes [D_2, \hat{M}_j].
  \label{eq:fluctuation}
\end{equation}
Both terms are generically nonzero: $[D, \hat{M}]$ is nonzero
whenever $\hat{M}$ has cross-shell matrix elements (coupling states
with different $n$, hence different $D$-eigenvalues). At $n_{\max} =
2$, four of 14 multiplier generators have nonzero commutators; at
$n_{\max} = 3$, 29 of 55.

The \emph{rotationally invariant} one-form is constructed using the
conjugate pairing~\cite{connes_book1994}:
\begin{equation}
  A = \sum_{(NLM)} \sum_{(N'L'M')}
  \bigl(\hat{M}_{NLM}^\dagger \otimes \hat{M}_{N'L'M'}^\dagger\bigr)
  \cdot [D_{\mathrm{total}}, \hat{M}_{NLM} \otimes \hat{M}_{N'L'M'}],
  \label{eq:A_full}
\end{equation}
summing \emph{independently} over all allowed multiplier labels on each
tensor factor in $\mathcal{O}_{n_{\max}}$. This is the full inner
fluctuation of the product algebra $\mathcal{A} \otimes \mathcal{A}$,
not only the diagonal ($N'L'M' = NLM$) terms; all numerical results in
this paper use this double sum. The dagger pairing ensures
$m$-conservation (total magnetic quantum number $M_1 + M_2$ is
preserved) and Hermiticity of $A$.

\begin{proposition}[Connected interaction]
\label{prop:connected}
The anti-commutator $\{D_{\mathrm{total}}, A\}$ contains a
non-factorizable (connected) two-body component. Define the connected
part via partial-trace subtraction:
\[
V_{\mathrm{conn}} = \{D,A\} - \frac{1}{d_2}\mathrm{Tr}_2\{D,A\}
\otimes I_2 - I_1 \otimes \frac{1}{d_1}\mathrm{Tr}_1\{D,A\}
+ \frac{\mathrm{Tr}\{D,A\}}{d_1 d_2} I_1 \otimes I_2.
\]
At $n_{\max} = 2$: $\|V_{\mathrm{conn}}\| / \|\{D,A\}\| = 76.7\%$.
At $n_{\max} = 3$: $\|V_{\mathrm{conn}}\| / \|\{D,A\}\| = 32.4\%$.
\end{proposition}

The interaction is genuine: the connected fraction is nonzero, and the
heat trace $\mathrm{Tr}\,e^{-t(D + \varepsilon A)^2}$ departs from
the factorized value monotonically with $\varepsilon$ (Table~\ref{tab:heat}).

\begin{table}[t]
\centering
\caption{Heat trace ratio $\mathrm{Tr}\,e^{-t(D+\varepsilon A)^2} /
  [\mathrm{Tr}\,e^{-tD_1^2} \cdot \mathrm{Tr}\,e^{-tD_2^2}]$ at
  $t = 0.5$ for the full-algebra gauge field.}
\label{tab:heat}
\begin{ruledtabular}
\begin{tabular}{ccc}
$\varepsilon$ & $n_{\max} = 2$ & $n_{\max} = 3$ \\
\hline
$0.0$ & $1.000000$ & $1.000000$ \\
$0.1$ & $0.995465$ & $0.838438$ \\
$0.5$ & $0.977815$ & $0.410165$ \\
$1.0$ & $0.956832$ & $0.185405$ \\
\end{tabular}
\end{ruledtabular}
\end{table}

%% =====================================================================
\section{Angular selection rules}
\label{sec:angular}
%% =====================================================================

\begin{theorem}[Angular structure of the connected interaction]
\label{thm:angular_structure}
The connected two-body interaction $V_{\mathrm{conn}}$ from the
rotationally invariant gauge field~\eqref{eq:A_full} satisfies:
\begin{enumerate}
\item[(i)] \textbf{$m$-conservation:} $m_1 + m_2 = m_1' + m_2'$ for
  all nonzero matrix elements. Verified at $100\%$ at $n_{\max} \in
  \{2, 3\}$.
\item[(ii)] \textbf{Gaunt compatibility:} $|\Delta l_1| = |\Delta
  l_2| = k$ for all nonzero $(l_1, l_2) \to (l_1', l_2')$ channels.
  Verified at $100\%$ at $n_{\max} \in \{2, 3\}$.
\item[(iii)] \textbf{Pure monopole:} $100\%$ of $\|V\|^2$ is in the
  $k = 0$ multipole channel at both cutoffs.
\item[(iv)] \textbf{Hermiticity and exchange symmetry} hold to
  machine precision ($< 10^{-17}$).
\end{enumerate}
\end{theorem}

\begin{proof}[Verification]
Direct computation from the $4$-index tensor $V_{\mathrm{conn}}[a, b,
c, d]$ extracted from the positive-chirality sector at $n_{\max} = 2$
($N = 5$, $25$ nonzero elements) and $n_{\max} = 3$ ($N = 14$, $196$
nonzero elements). The $m$-conservation is enforced by the $M^\dagger
\cdot [D, M]$ conjugate pairing, which ensures the total $M$ quantum
number of the gauge field is zero. The Gaunt compatibility follows
from the $3$-$Y$ selection rules of the multiplier algebra.
\end{proof}

\begin{remark}
The $k = 0$ monopole channel corresponds to the $F^0$ Slater integral
in the Neumann expansion of $1/r_{12}$
(Paper~12~\cite{paper12}). Higher multipoles $k = 1, 2, \ldots$ would
appear at larger $n_{\max}$ or in the full spinor sector. The
\emph{absence} of the $k = 1$ dipole is correct for identical
particles by parity.
\end{remark}

%% =====================================================================
\section{Radial comparison with Coulomb}
\label{sec:radial}
%% =====================================================================

\begin{proposition}[Radial weights do not match Coulomb]
\label{prop:radial_negative}
The matrix element ratios $V_{\mathrm{conn}}[a,b,c,d] /
\langle ab | r_{12}^{-1} | cd \rangle$ are not constant across the
nonzero elements. At $n_{\max} = 2$: coefficient of variation $\mathrm{CV}
= 16.5$, Pearson correlation $r = 0.58$, sign agreement $68\%$. At
$n_{\max} = 3$: $\mathrm{CV} = 20.3$, $r = 0.41$, sign agreement
$54\%$.
\end{proposition}

The spectral-action interaction is a strict \emph{subset} of the
Coulomb matrix elements: all nonzero spectral elements also appear in
the Coulomb matrix ($0$ spectral-only entries), but the Coulomb matrix
has many additional nonzero elements ($240$ at $n_{\max} = 2$,
$15\,168$ at $n_{\max} = 3$) that are zero in the spectral
interaction. The Pearson correlation \emph{decreases} with $n_{\max}$,
indicating this is structural rather than a finite-size artifact.

\begin{remark}
The $1/r_{12}$ Coulomb kernel in the Fock-projected basis arises from
the chordal distance identity under stereographic projection
(Paper~7~\cite{paper7}). This is a \emph{metric} property of $S^3$,
not a spectral-action property. The spectral action
$\mathrm{Tr}\,f(D^2/\Lambda^2)$ computes heat-kernel coefficients
(Seeley--DeWitt), while $1/r_{12}$ is a Green's function kernel ---
fundamentally different objects at the mathematical level.
\end{remark}

%% =====================================================================
\section{The algebra/metric boundary}
\label{sec:boundary}
%% =====================================================================

The results of Sections~\ref{sec:angular}--\ref{sec:radial} precisely
delineate the universal/Coulomb partition of
Paper~31~\cite{paper31} at the two-body level:

\begin{itemize}
\item The \textbf{algebra} $\mathcal{A}$ (truncated operator system,
  $3$-$Y$ multipliers) determines the angular selection rules: Gaunt
  compatibility, $m$-conservation, and the multipole hierarchy.
  These are \emph{universal} --- they hold for any central potential,
  as proved by the angular sparsity theorem of
  Paper~22~\cite{paper22}.

\item The \textbf{Dirac operator} $D$ (Camporesi--Higuchi eigenvalues)
  determines the radial coupling strengths. The gauged spectral action
  generates an interaction whose radial weights depend on the specific
  spectrum of $D$, not on the $1/r$ Coulomb potential. The Coulomb
  kernel is additional input from the Fock projection's metric
  structure.
\end{itemize}

This decomposition explains the ``multi-focal composition wall''
documented across six observables in the GeoVac
programme~\cite{paper34}: the wall sits at the $\mathcal{A}/D$
boundary. The framework couples discrete labels (angular quantum
numbers) via the algebra's inner fluctuations, but the spatial
coupling strength (radial integrals) requires the metric --- the Fock
projection that maps the graph onto $S^3$, and with it, the $1/r$
potential.

\begin{corollary}
For two particles on $\mathcal{T}_{S^3} \otimes \mathcal{T}_{S^3}$,
the inter-particle interaction splits as
\begin{equation}
  V_{12} = \underbrace{\text{angular selection
      rules}}_{\text{from } \mathcal{A}} \times
  \underbrace{\text{radial coupling
      strengths}}_{\text{from } D / \text{metric}}.
  \label{eq:AD_split}
\end{equation}
The first factor is framework-native (derived from inner
fluctuations). The second factor is calibration data (the Slater
integrals $R^k$ of Paper~12~\cite{paper12}).
\end{corollary}

%% =====================================================================
\section{Implications}
\label{sec:implications}
%% =====================================================================

\subsection{For atomic and molecular structure}

The result vindicates the composed-geometry architecture of
Paper~17~\cite{paper17}: the Gaunt selection rules that give the
composed Pauli scaling $O(Q^{2.5})$ of Paper~14~\cite{paper14} are
not ad hoc --- they emerge from the tensor-product algebra's inner
fluctuations. The radial Slater integrals that complete the
Hamiltonian are the irreducible calibration content, consistent with
Paper~18's exchange-constant taxonomy~\cite{paper18} (calibration
tier).

\subsection{Connection to double factorization and tensor
hypercontraction}
\label{subsec:df_thc_connection}

The two-electron operator decomposition $V_{ee} = \sum_P L_P \otimes
L_P$ used in double factorization~\cite{vonBurg2021} and tensor
hypercontraction~\cite{Lee2021} is literally an element of
$\mathcal{A} \otimes \mathcal{A}$ in the tensor-product spectral
triple constructed here. In the GeoVac basis, this decomposition is
supplied analytically by the multipole expansion
\begin{equation*}
  \frac{1}{r_{12}} = \sum_{L,M}\frac{4\pi}{2L{+}1}\,
  \frac{r_<^L}{r_>^{L{+}1}}\,Y_{LM}(\Omega_1)\,Y_{LM}^{*}(\Omega_2),
\end{equation*}
with $L_P$ proportional to $Y_{LM}$ and rank equal to the number of
distinct radial-density channels $\rho_\alpha(r) = R_{n_a l_a}(r)\,
R_{n_b l_b}(r)$ at each angular momentum $L$.

Numerical double factorization applied to a GeoVac composed ERI
tensor reproduces this analytic decomposition bit-exactly at the
production basis: the relative reconstruction error
$\|G - V K V^\top\|_F / \|G\|_F$ between the empirical ERI tensor and
its multipole-VKV form is $10^{-16}$ at $n_{\max} = 2$ and $10^{-14}$
at $n_{\max} = 3$, i.e.\ machine precision. Every singular vector of
the reshaped two-electron tensor lives strictly inside the
multipole-channel subspace at every basis size tested
($n_{\max} \in \{2, 3, 4\}$), with total subspace overlap unity to
machine precision. At the production basis ($n_{\max} = 2$, valence
$\{1s, 2s, 2p\}$), the double-factorization rank equals the analytic
distinct-radial-density count exactly --- seven, distributed as four
densities at $L = 0$, two at $L = 1$, one at $L = 2$.

At larger basis ($n_{\max} \geq 3$), the radial-integral matrix
$R^L(\rho_\alpha; \rho_\beta) = \langle\rho_\alpha|r_<^L/r_>^{L+1}|
\rho_\beta\rangle$ develops a graded singular-value spectrum that
decays over many orders of magnitude rather than terminating at a
clean cutoff. This is the well-known
Laguerre-basis ill-conditioning of hydrogenic ERIs --- the same
phenomenon that motivates resolution-of-identity and Cholesky
methods in standard atomic and molecular structure. Radial densities
with similar effective exponents $(1/n_a + 1/n_b)/2$ produce nearly
proportional contributions against the Coulomb kernel, so the
apparent ``rank'' of $R^L$ depends on threshold. The additional compression that DF finds beyond the bare
analytic multipole count is therefore not a new algebraic structure
but a numerical realization of the same spectral content, exposed
through SVD ordering. Operationally: double factorization on a
GeoVac Hamiltonian is the multipole expansion accessed by SVD rather
than by closed form. Detailed verification in
\texttt{debug/sprint\_df\_multipole\_lift\_memo.md}.

\subsection{For nuclear structure}

The same $\mathcal{A}/D$ split applies to the nucleon--nucleon
interaction. The Minnesota potential's angular structure (central
force, spin-dependent channels) is algebra-side; the radial Gaussian
form factors ($V_R e^{-\kappa_R r^2}$, etc.)\ are $D$-side
calibration. This explains why the nuclear model requires imported
potentials: the framework derives \emph{which} nuclear channels couple
but not the NN force strength.

\subsection{For gravity}

The gravity sector (Paper~51~\cite{paper51}) is special: the spectral
action directly gives metric properties (Einstein--Hilbert action,
Newton's constant, cosmological constant) because gravity \emph{is}
the metric. The two-body Coulomb interaction is a Green's function
kernel, not a heat-kernel coefficient, so it lives on the other side
of the $\mathcal{A}/D$ boundary. This structural difference explains
why gravity results are bit-exact while inter-particle interactions
require imported potentials.

%% =====================================================================
\section{Discussion}
\label{sec:discussion}
%% =====================================================================

The tensor-product spectral action generates genuine two-body physics:
a non-factorizable interaction with the correct angular selection
rules. That the radial weights do not match Coulomb is not a failure
of the framework but a structural theorem about \emph{what spectral
actions can compute}: heat-kernel coefficients (metric curvature, not
Green's functions). The Coulomb potential is the Green's function of
the Laplacian, $\nabla^2 G = -4\pi\delta^{(3)}$; it requires solving
a differential equation on the manifold, not tracing over the
spectrum.

The practical consequence is that inter-particle radial integrals
(Slater $R^k$, Minnesota form factors, Clementi--Raimondi screening)
are irreducible calibration data in the GeoVac framework. The angular
selection rules that govern which integrals are nonzero are
framework-native. This is a more precise statement of the multi-focal
composition wall: it is not a missing tool but the structural boundary
between what the algebra determines and what the metric determines in
a spectral triple.

\emph{Open question.} Whether the Green's function kernel can be
recovered from the spectral triple by a construction other than the
spectral action --- for example, via the resolvent $(D - z)^{-1}$ or
the Connes distance function --- remains an open problem. The
resolvent naturally gives the propagator $\langle x | (D-z)^{-1} |
y\rangle$, which in the continuum limit reproduces the Green's
function. A resolvent-based derivation of $1/r_{12}$ from the
tensor-product spectral triple would close the multi-focal wall.

\begin{acknowledgments}
This work was performed using AI-augmented research methods.
Computational implementation and analysis were carried out
collaboratively with Anthropic Claude.
\end{acknowledgments}

\begin{thebibliography}{30}

\bibitem{paper7} J.~Loutey, ``The Dimensionless Vacuum: Recovering the
  Schr\"odinger Equation from Scale-Invariant Graph Topology,'' GeoVac Paper~7
  (2025).

\bibitem{paper12} J.~Loutey, ``Algebraic Two-Electron Integrals on the
  Prolate Spheroidal Lattice,'' GeoVac Paper~12 (2025).

\bibitem{paper14} J.~Loutey, ``Structurally Sparse Qubit
  Hamiltonians from Spectral Graph Theory,'' GeoVac Paper~14 (2025).

\bibitem{paper17} J.~Loutey, ``Composed Natural Geometries for
  Core-Valence Diatomics: LiH and BeH$^+$ from Discrete Graph Topology,''
  GeoVac Paper~17 (2025).

\bibitem{paper18} J.~Loutey, ``Spectral-Geometric Exchange
  Constants: Projecting Discrete Spectra onto Continuous Manifolds,''
  GeoVac Paper~18 (2026).

\bibitem{paper22} J.~Loutey, ``Potential-Independent Angular Sparsity
  for Qubit Hamiltonians in Spherical Fermion Systems,'' GeoVac
  Paper~22 (2026).

\bibitem{paper31} J.~Loutey, ``The Universal/Coulomb Partition of the
  GeoVac Spectral Triple,'' GeoVac Paper~31 (2026).

\bibitem{paper32} J.~Loutey, ``The GeoVac Spectral Triple: A Synthesis
  Paper Locking the Framework into the Marcolli--van~Suijlekom
  Gauge-Network Lineage,'' GeoVac Paper~32 (2026).

\bibitem{paper34} J.~Loutey, ``GeoVac as a Two-Layer Framework:
  Projection Taxonomy and the Empirical Matches Catalog,'' GeoVac Paper~34 (2026).

\bibitem{paper38} J.~Loutey, ``State-space Gromov--Hausdorff convergence
  of truncated Camporesi--Higuchi spectral triples on $S^3$,''
  GeoVac Paper~38 (2026).

\bibitem{paper51} J.~Loutey, ``Gravity from the GeoVac spectral action:
  continuum results and the discrete-substrate program,'' GeoVac Paper~51 (2026).

\bibitem{camporesi_higuchi1996} R.~Camporesi and A.~Higuchi, ``On the
  eigenfunctions of the Dirac operator on spheres and real hyperbolic
  spaces,'' J.~Geom.\ Phys.\ \textbf{20}, 1 (1996).

\bibitem{connes_book1994} A.~Connes, \textit{Noncommutative Geometry}
  (Academic Press, 1994).

\bibitem{connes_vs2021} A.~Connes and W.~D.\ van Suijlekom,
  ``Spectral truncations in noncommutative geometry and operator
  systems,'' Commun.\ Math.\ Phys.\ \textbf{383}, 2021 (2021).

\bibitem{marcolli_vs2014} M.~Marcolli and W.~D.\ van Suijlekom,
  ``Gauge networks in noncommutative geometry,'' J.~Geom.\ Phys.\
  \textbf{75}, 71 (2014).

\bibitem{CLAUDE_wall} J.~Loutey, ``The Multi-Focal Composition Wall:
  Five Independent Observables, Same Structural Reason,'' GeoVac
  project documentation (2026).

\bibitem{vonBurg2021} V.~von Burg, G.~H.~Low, T.~H{\"a}ner,
  D.~S.~Steiger, M.~Reiher, M.~Roetteler, and M.~Troyer, ``Quantum
  computing enhanced computational catalysis,'' Phys.\ Rev.\ Research
  \textbf{3}, 033055 (2021).

\bibitem{Lee2021} J.~Lee, D.~W.~Berry, C.~Gidney, W.~J.~Huggins,
  J.~R.~McClean, N.~Wiebe, and R.~Babbush, ``Even More Efficient
  Quantum Computations of Chemistry Through Tensor Hypercontraction,''
  PRX Quantum \textbf{2}, 030305 (2021).

\end{thebibliography}

\end{document}
